As conclusões desta replicação se dividem em duas. Em primeiro lugar, o mecanismo se
transfere, uma vez que a AGOP recupera direções lineares de conceito mesmo na escala de
8{,}73M de parâmetros, alcançando um $|\rho|$ de validação de até 0{,}978, e essas direções são
causais quando escritas de volta no fluxo residual. Em segundo lugar, a capacidade que o
\textit{steering} deveria trazer não aparece, dado que nesta escala a intervenção move o
vocabulário que o modelo busca muito antes de produzir uma frase coerente carregando o
conceito, de forma que obtemos um mecanismo sem o comportamento que ele deveria sustentar.

A afirmação comparativa, por sua vez, jamais recebeu um teste justo. A RFM de fato bate a
PCA de forma limpa, com um $|\rho|$ de validação de 0{,}918 contra 0{,}154 em \verb|valence| e
uma AUROC de \textit{monitoring} de 1{,}000 contra 0{,}879 no agregado, e ainda assim RFM, DiffMeans e
LogReg se sentam juntas no teto da AUROC, com a LogReg inclusive superando a RFM no
$|\rho|$ de validação nos três conceitos. Note que
isso diz que os três métodos são indistinguíveis \emph{por estas métricas e sobre estes
conceitos}, e não diz que a AGOP recupera a mesma direção que uma subtração de duas médias
de classe, o que é uma afirmação à parte, testável pelo alinhamento por bloco
$\cos\left(v_{\mathrm{RFM}}, v_{\mathrm{DiffMeans}}\right)$, e que não calculamos.

Uma consequência metodológica deste trabalho é portável para qualquer escala, e ela não é
um resultado sobre o GuppyLM. Um coeficiente de \textit{steering} deve ser reportado
relativamente à escala das ativações, na forma
$\varepsilon_\ell = r \cdot \operatorname*{median}_{i} \left\| a_\ell^{(i)} \right\|_2$,
uma vez que as normas por bloco aqui medidas percorrem de 11{,}9 a 30{,}9 enquanto os valores
absolutos de Beaglehole et al. são calibrados para a escala de ativações do Llama, de forma
que é o coeficiente relativo $r$, e não $\varepsilon$, que se transporta entre modelos.
Note que uma replicação conduzida em qualquer outra escala herda esse problema antes de
herdar qualquer outro.

Para sustentar as demais afirmações substantivas que Beaglehole et al.
(\cite{beaglehole2026toward}) fazem, a saber, o escalonamento da \textit{steerability} com o
tamanho do modelo, a vantagem da RFM sobre conceitos difíceis, e os referenciais reais de
\textit{monitoring} como FAVABENCH, HaluEval e ToxicChat, necessitamos de uma infraestrutura que
um único \textit{checkpoint} do GuppyLM não é capaz de fornecer. Tomaremos, portanto, uma
escada de modelos Pythia como plano subsequente, analisando \textit{checkpoints} ao longo de
diversos tamanhos de modelo de modo a ler a \textit{steerability} diretamente contra a escala,
ao invés de inferi-la. A questão principal ali é se a escada passa a produzir texto fluente
carregando o conceito ao invés de deriva de vocabulário, e em que tamanho ela passa a
fazê-lo.

Em suma, a maquinaria de Beaglehole et al. parece não requerer um modelo grande para
funcionar, mas apenas para valer a pena, e é a distância entre essas duas afirmações que um
único \textit{checkpoint} minúsculo é capaz de medir. Discutimos aqui um modelo, uma
família de conceitos dentro do domínio e dois juízes imperfeitos, e muito mais ainda há a
ser discutido, particularmente o comportamento do mesmo \textit{pipeline} sobre conceitos
que um modelo deste tamanho não enuncie já lexicalmente em sua própria saída.
